\clearpage
\SourceLeaf{404}
\begin{center}\large [GENERALIZATION OF LEGENDRE'S THEOREM.]\end{center}
\begin{center}[I.]\end{center}
\begin{center}\textsc{Gauss} to \textsc{Olbers}. Göttingen, 9 October 1825.\end{center}
\indent\mbox{. . . . . . . . . .} In recent days I have begun, with an eye to my future work on Higher Geodesy, to put a (very) small part of the material on curved surfaces into some order in my mind. Yet I am becoming convinced that, given the distinctive character of my whole treatment, I have to \textls{start very far back} for the sake of coherence. I must even set out beforehand my view of the radii of curvature of \textls{plane} curves. This has made me wonder whether it would be wiser first to work out separately a part of this theory that is purely geometrical (in analytical form) and presents new and familiar material together in a new form, perhaps to detach it from the work and include it in our Commentationen as one or two papers. For now, however, I can leave the form of publication open and continue putting it on paper. \mbox{. . .}
